\section{Coloring a nonsingular measured graph}
\label{sec:coloring}

The polynomial calculations of Sections~\ref{sec:polynomials} and
\ref{sec:invariant} will now produce a measurable coloring on a probability
extension of a nonsingular measured graph.  The conclusion concerns the
operator norm on a prescribed finite ball, with an arbitrarily small
probability of failure.  We will also preserve the uniform distribution of
each color conditional on the original base point.  This latter condition
will be needed when the coloring is used inside a von Neumann algebra.

There are two steps in passing from the preceding invariant calculations to
this conclusion.  We first construct costs on the nonsingular probability
space by adding a real height coordinate and then removing upper and lower
weight cutoffs.  We then assign colors at a small random fraction of the
remaining vertices at each step.  An estimate for the integrated difference
of the limiting costs, on the set where two colorings agree on a large ball,
will control the interaction between these simultaneous assignments.

\subsection{The probability extension and the coloring theorem}

Let $(X,\mu)$ be a standard probability space, and let $\mathcal R$ be a
countable nonsingular measured equivalence relation on $X$.  We use Borel
representatives after discarding null sets when necessary.  Write the
Radon--Nikodym cocycle in the orientation
\begin{equation}
 \mu(\gamma S)
   =\int_S\Delta(\gamma v,v)\,d\mu(v)
 \label{eq:coloring-cocycle}
\end{equation}
for measurable partial isomorphisms $\gamma$ with graph in $\mathcal R$ and
measurable sets $S$ in their domains.  Versions can be chosen off a saturated
null set so that the cocycle identity holds.  For example, one first uses
countably many partial isomorphisms covering $\mathcal R$; the chain rule
gives the identity for their compositions, and derivatives agree on sets
where two such maps coincide.  Throughout this Section, null sets arising
along the countable relation may be discarded together with their
saturations, which are null by nonsingularity.

Let $G$ be an undirected measurable graph contained in $\mathcal R$, without
loops, and with degree at most $D_0<\infty$.  Its radius-$m$ ball at $v$ is
denoted by $B_G(v,m)$.  Suppose that
\begin{equation}
 K^{-1}\leq\Delta(u,v)\leq K
 \qquad\text{on every edge of }G,
 \qquad K\geq1.
 \label{eq:coloring-edge-ratios}
\end{equation}
We consider a measurable field of Hermitian matrices $T$ on the orbit
spaces $\ell^2([v]_{\mathcal R})$, with norm at most $1$, zero diagonal, and
entries supported on $G$.  Upon changing the root within an orbit, the
matrices are identified up to conjugacy by a diagonal unitary.  All
quantities used below, including principal norms and the polynomial costs,
are invariant under these conjugacies.

The \emph{relative Bernoulli extension} of $(X,\mu,\mathcal R)$ is the space
\[
 \widetilde X
   =\{(v,\omega):v\in X,\ \omega\in[0,1]^{[v]_{\mathcal R}}\},
 \qquad \pi(v,\omega)=v,
\]
equipped with the probability measure $\widetilde\mu$ whose conditional
measure over $v$ is the product of Lebesgue measures, denoted by $\kappa_v$.
The equivalence relation on this extension identifies $(v,\omega)$ with
$(u,\omega)$ whenever $(u,v)\in\mathcal R$, regarding $\omega$ as one
function on the common class.  Each extended orbit thus projects bijectively
onto a base orbit.  The graph and matrices are pulled back along this
identification.

These are standard measurable spaces and fields: enumerate the base
relation by countably many partial isomorphisms, and use one coordinate for
each distinct site in a class.  Changing the root reindexes a product
family and preserves its conditional product measure.  Consequently the
lift of a base partial isomorphism has the same Radon--Nikodym derivative
as that base map.  These lifts enumerate the extended relation.  Cutting
their domains by extension-measurable sets and taking countable disjoint
unions therefore shows, by conditional integration, that the cocycle on
the extension is the pullback of $\Delta$.  A single Lebesgue coordinate
at a site will be used to encode countably many independent random
coordinates whenever needed.

For an integer $r\geq1$ and a coloring
$c:\widetilde X\to\{1,\ldots,r\}$, its value at a site $u$
in the orbit of $(v,\omega)$ means $c(u,\omega)$.  On a finite principal
set $I$ in this orbit, let
\[
 T_{I,c}
   =\sum_{i=1}^r 1_{\{c=i\}\cap I}\,T|_I\,
                         1_{\{c=i\}\cap I}
\]
be the matrix obtained by retaining entries joining vertices of the same
color.

\begin{theorem}[Bernoulli coloring]
\label{thm:bernoulli-coloring}
Let $(X,\mu,\mathcal R)$, $G$, $\Delta$, and $T$ satisfy the assumptions
above: $\mathcal R$ is countable and nonsingular on a standard probability
space, $G$ is undirected, has no loops, has degree at most $D_0<\infty$, and satisfies
Equation~\eqref{eq:coloring-edge-ratios}, and $T$ is a measurable
Hermitian contraction with zero diagonal and support on $G$, allowing
diagonal unitary conjugacy under a change of root.  Fix an integer $r\geq1$
and a real number $t>3/\sqrt r$, and put $\beta=80$.

For every integer $m\geq1$ and every $\alpha,\eta>0$, there is a measurable
coloring $c:\widetilde X\to\{1,\ldots,r\}$ on the relative Bernoulli
extension such that
\begin{equation}
 \kappa_v\{\omega:c(v,\omega)=i\}=\frac1r
 \quad\text{for every }i\in\{1,\ldots,r\}
 \quad\text{and almost every }v\in X,
 \label{eq:coloring-marginal}
\end{equation}
and
\begin{equation}
 \widetilde\mu\left\{(v,\omega):
       \|T_{B_G(v,m),c}\|\leq2\beta t+\alpha\right\}>1-\eta.
 \label{eq:coloring-ball-conclusion}
\end{equation}
\end{theorem}

We prove the theorem after constructing the cost that will control the
assignments.  Until the final proof, the probability base may already
include some Bernoulli coordinates.  We continue to denote this current
base by $(X,\mu)$ and its relation by $\mathcal R$; the cocycle and all graph
and matrix bounds remain as above.

\subsection{Homogeneous costs from an invariant lift}

A partial coloring is a measurable function
$c:X\to\{0,1,\ldots,r\}$, with $0$ indicating an unassigned site.  Recall
the finite-site data of Section~\ref{sec:polynomials}: the diagonal shift
at a site is
\[
 \tau_v=
 \begin{cases}
 t,&c(v)=0,\\
 \beta t,&c(v)>0,
 \end{cases}
\]
and the allowed polynomial choices are both signs of all $r$ colors at an
unassigned site, or both signs of its assigned color otherwise.  Choices
are uniform among the allowed pairs and independent across sites.
With positive weights $W_v$, the outcome
matrix on a choice block is the corresponding principal restriction of
\[
 \diag(\sqrt W)(\pm T)\diag(\sqrt W)-\diag(\tau W).
\]
For the expected characteristic polynomial $P$ and its single-site
deletion $Q=P_{-v}$, we use the notation of that Section:
\[
 F_v=P/Q,\qquad G_v=1/F_v,
 \qquad m_v=\int x_+\,\rho_v(dx),
 \qquad d_v=s(P)-s(Q),
\]
where $\rho_v$ is the rooted probability measure with Cauchy transform
$G_v$, $x_+=\max(x,0)$, and $s(P)$ is the sum of the positive parts of the
roots, counted with multiplicity.  At an unassigned site the superscript
$i$ denotes permanent assignment to color $i$, including the increase in
the shift.  The quotient $F_{v,i}^0$ denotes conditioning on color $i$
without increasing that shift, whereas
\[
 F_{v,i}=F_{v,i}^0+(\beta-1)tW_v
\]
includes the increase.  We set the virtual fields $d_v^{i}$ to zero on
already assigned sites.  On bounded invariant measured graphs the same
notation refers to the limits supplied by Lemma~\ref{lem:bounded-costs}.

The original measure need not be invariant, so the integral change formula
cannot yet be applied there.  We use the Maharam invariant lift
\cite{Mah64}; for the multiplicative-coordinate formula, see
\cite[Section~2, before Theorem~2.1]{RoySamorodnitsky06}.
Here we use logarithmic height, with multiplicative coordinate $y=e^s$;
the cocycle verification and subsequent homogeneous-cost construction are
given below.  Adjoin a height $s\in\mathbb R$ and set
\begin{equation}
 d\nu(v,s)=d\mu(v)e^s\,ds,
 \qquad W(v,s)=e^{-s}.
 \label{eq:coloring-lift}
\end{equation}
In the lifted relation, the site corresponding to $u$ at a root $(v,s)$ is
\begin{equation}
 \bigl(u,s-\log\Delta(u,v)\bigr).
 \label{eq:coloring-height-change}
\end{equation}
The measure $\nu$ is invariant.  Indeed, the factor $\Delta(u,v)$ from
Equation~\eqref{eq:coloring-cocycle} cancels the factor
$e^{-\log\Delta(u,v)}$ from the height density.  The destination weight is
$\Delta(u,v)W(v,s)$, so neighboring weights have ratios in $[K^{-1},K]$.
The coloring and matrix data are independent of height.

For $0<l<H<\infty$, retain only the sites with $l\leq W\leq H$, using the
restricted relation, graph, and principal matrix data.  The retained
measure is finite, since
\[
 \nu\{l\leq W\leq H\}=l^{-1}-H^{-1}.
\]
All results of Section~\ref{sec:invariant} therefore apply.  Denote the
resulting cost fields by $m^{l,H},d^{l,H},(d^i)^{l,H}$.  The next lemma
removes both cutoffs, while retaining control uniform enough for later
changes of coloring.

\begin{lemma}[Homogeneous costs]
\label{lem:homogeneous-costs}
Fix $D_0,K,r,t$ and $\beta=80$, and let the probability space, relation,
graph, matrix data, and partial coloring satisfy the assumptions above.
For each field $f$ in the finite list $m,d,d^1,\ldots,d^r$, there is a
bounded measurable function $\bar f$ on $X$ such that, on every fixed band
\[
 \mathcal U_{b_1,b_2}=\{(v,s):b_1\leq W(v,s)\leq b_2\},
 \qquad 0<b_1<b_2<\infty,
\]
one has
\begin{equation}
 \frac{f^{l,H}}{W}\longrightarrow\bar f(v)
 \quad\text{in }L^1(\mathcal U_{b_1,b_2},\nu)
 \quad\text{jointly as }l\downarrow0, H\uparrow\infty.
 \label{eq:coloring-band-limit}
\end{equation}
The errors tend to zero uniformly over all such probability bases and
configurations with the fixed parameters and bounds.  There is a constant
$C_0$, depending only on those parameters, such that all the normalized
cutoff fields on retained sites and all their limits have absolute value
at most $C_0$.
The limit fields are independent of height and satisfy
\begin{equation}
 \min_{1\leq i\leq r}\bar d^{i}-\bar d
      \leq24\bar m
 \quad\text{almost everywhere on }\{c=0\}.
 \label{eq:coloring-limit-assignment}
\end{equation}
\end{lemma}

\begin{proof}
The two cutoffs use different controls. The set of weights above $H$ has
invariant measure $H^{-1}$, so the rank comparison controls the upper
end. With the upper cutoff fixed, the edge-ratio bound makes sufficiently
small weights invisible in any chosen finite neighborhood of a fixed band.
We first prove uniform convergence on the band $1\leq W\leq e$.  The same
argument applies to any fixed band bounded away from zero.  Keep the lower
cutoff $l$ fixed and compare upper cutoffs $H_2\geq H_1\geq e$.  Passing
from $H_2$ to $H_1$ deletes a set of invariant measure at most $1/H_1$.
Lemma~\ref{lem:rank-comparison} gives, on the appropriate sites of the band,
integrated absolute differences of order $1/H_1$ for
\[
 G_v(z)
 \quad\text{and}\quad
 J_i(v,z)=\bigl(F_{v,i}^0(z)-z-a_v\bigr)G_v(z),
 \qquad a_v=\tau_vW_v.
\]
For $J_i$ the targets are the unassigned band sites.  The constants are
uniform in $l,H_1,H_2$, the configuration, and $z$ on each compact subset
of the open upper half-plane.

The second-moment estimate of Lemma~\ref{lem:rooted-locality} is
\begin{equation}
 \int x^2\,\rho_v(dx)
       \leq\bigl((\beta t)^2+K\bigr)W_v^2.
 \label{eq:coloring-second-moment}
\end{equation}
It also holds for each of the altered site choices used here.  In
particular, these second moments are uniformly bounded on the band.  They
give a uniform positive lower bound for $|G_v(z)|$ on each compact subset
of the upper half-plane: a fixed bounded interval carries a uniformly
positive amount of the probability measure, and
\[
 -\operatorname{Im}G_v(z)
   =(\operatorname{Im}z)\int |z-x|^{-2}\,\rho_v(dx).
\]
The residue bound in Lemmas~\ref{lem:polynomial-quotients} and
\ref{lem:rooted-locality} also bounds $J_i$ uniformly on these compact
sets.  Division is thus controlled in the identities
\[
 F_v=1/G_v,
 \qquad F_{v,i}^0=z+a_v+J_i/G_v,
 \qquad F_{v,i}=F_{v,i}^0+(\beta-1)tW_v.
\]
The integrated differences of these quotients consequently tend to zero
as $H_1\to\infty$, uniformly in the lower cutoff.

We next recover the costs from these transform comparisons.  Probability
measures with the uniform second-moment bound form a weakly compact
family, and integration of $x_+$ is continuous on this family by uniform
integrability.  Cauchy transforms on a countable dense subset of the upper
half-plane separate probability measures.  Compactness therefore implies
that, for any prescribed accuracy of $m_v$, comparisons of the transforms
at finitely many arguments to sufficiently small accuracy suffice.
The integrated transform estimates, together with the uniform bound on
$m_v$, give the desired integrated comparison for $m$.

For deletion costs we use the argument of the quotient.  In a finite
restriction, let $n_P(x)$ and $n_Q(x)$ count roots strictly greater than
$x$, with multiplicity.  Interlacing and the quotient properties give,
for $y>0$,
\begin{equation}
 \arg F(x+iy)
   =\int_{\mathbb R}\frac{y}{(x-u)^2+y^2}
                     (n_P-n_Q)(u)\,du,
 \qquad 0\leq\arg F\leq\pi.
 \label{eq:coloring-argument}
\end{equation}
This follows by writing the arguments of the individual root factors;
the branch is fixed by $\operatorname{Im}F>0$.  The spectral-shift function
$n_P-n_Q$ takes values between zero and one.  Lemma~\ref{lem:positive-cost}
and Equation~\eqref{eq:coloring-second-moment} imply that
\begin{equation}
 \int_B^\infty(n_P-n_Q)(x)\,dx=O(B^{-1})
 \quad(B\to\infty),
 \label{eq:coloring-shift-tail}
\end{equation}
uniformly on the band, including for the shifted virtual options.  Indeed,
the bound in that lemma is
\[
 2\int(x-B)_+\,\rho_v(dx)+(a_v+B)\rho_v\{x>B\},
\]
whose two terms are controlled by the second moment and the band bound on
$a_v$.

On $[0,B]$, the integral of the spectral-shift function is approximated by
\begin{equation}
 \frac1\pi\int_0^B\arg F(x+iy)\,dx
 \quad\text{as }y\downarrow0,
 \label{eq:coloring-shift-poisson}
\end{equation}
uniformly over these finite calculations.  To see uniformity, interchange
the integrals in Equation~\eqref{eq:coloring-argument}; the error is at most
the $L^1(\mathbb R)$ error in the Poisson approximation to $1_{[0,B]}$,
because $0\leq n_P-n_Q\leq1$. More precisely, put
\[
 p_y(u)=\frac{y}{\pi(u^2+y^2)},\qquad
 \delta_B(y)=\|p_y*1_{[0,B]}-1_{[0,B]}\|_{L^1(\mathbb R)}.
\]
For fixed $B$, one has $\delta_B(y)\to0$ as $y\downarrow0$. Combining
the smoothing error with Equation~\eqref{eq:coloring-shift-tail} gives,
uniformly for the finite rooted calculations on the band,
\begin{equation}\label{eq:coloring-cut-cost-recovery}
 \left|d_v-\frac1\pi\int_0^B\arg F_v(x+iy)\,dx\right|
 \leq\frac{C}{B}+\delta_B(y),\qquad B\geq1.
\end{equation}
The same estimate holds for each shifted virtual option. For fixed $B,y$,
convergence of the costs and locally uniform convergence of the quotients
in Lemma~\ref{lem:bounded-costs} pass this estimate to every present band
vertex of a bounded infinite cut, and to present unassigned band vertices for
the virtual options. The quotient has imaginary part at least $y$, so its argument
is continuous throughout this passage. We use this estimate in terms of
the cost and quotient on the infinite cut. For fixed $B,y$, the arguments in its
integral are controlled by the integrated quotient comparisons already
proved. First choose $B$ large, then $y$ small for that $B$, and finally
the upper cutoff $H_1$ large for the fixed segment $[0,B]+iy$. This proves
the high-cutoff comparison for $d$ and every $d^i$.

It remains to remove the lower cutoff.  For a fixed finite upper cutoff,
Lemma~\ref{lem:bounded-costs} approximates the costs uniformly by bounded
neighborhood calculations.  A radius-$R$ neighborhood of a root in the
fixed band cannot contain a vertex of weight below $K^{-R}$.  Thus a
sufficiently small lower cutoff is invisible to any one of these fixed
neighborhood calculations.  This proves uniform convergence as
$l\downarrow0$ with that upper cutoff fixed.  For joint convergence, first
compare both large upper cutoffs to one sufficiently large fixed upper
cutoff by the preceding estimates, and then compare the lower cutoffs
there.  This gives uniform joint $L^1$-Cauchy estimates.  Division by $W$
does not change convergence on a fixed band.

The uniform bound asserted in the lemma follows directly from
Equation~\eqref{eq:coloring-second-moment} and
Lemma~\ref{lem:positive-cost}: $m_v/W_v$ is uniformly bounded, and
$d_v\leq2m_v+a_v$.  The same bounds hold for all virtual options.  In
particular they are independent of the band, the cutoff, and the
configuration.

We have constructed limits on every fixed band.  To identify their height
dependence, translate height by a real number $b$.  This scales every
weight by $e^{-b}$ and changes the cutoffs correspondingly.  All outcome
matrices and all costs scale by the same positive factor.  Thus
\[
 \frac{f^{l,H}(v,s+b)}{W(v,s+b)}
   =\frac{f^{e^b l,e^b H}(v,s)}{W(v,s)}.
\]
Joint convergence makes the limiting normalized field invariant almost
everywhere under every rational height translation, simultaneously.
A measurable function on $\mathbb R$ invariant almost everywhere under
all rational translations is constant almost everywhere; this follows,
for example, by convolution.  Applying this on almost every base fiber
gives the functions $\bar f(v)$ and their asserted independence of height.
Their measurability also follows from the band limits, for instance by
integrating a representative over a fixed height interval.

Finally Lemma~\ref{lem:assignment} gives the single-site inequality on
finite restrictions.  Lemma~\ref{lem:bounded-costs} passes it to each
bounded cut.  Joint band convergence then permits a pointwise convergent
subsequence for the finite list of normalized fields.  Dividing by the
positive root weight and taking this limit yields
Equation~\eqref{eq:coloring-limit-assignment}.
\end{proof}

We record two compatibility properties of this construction for later use.
Suppose $(\widehat X,\widehat\mu,\widehat{\mathcal R})$ is a standard
nonsingular extension with a measurable map $\pi:\widehat X\to X$ such that
$\pi_*\widehat\mu=\mu$, each relation class projects bijectively onto an
original class, and its Radon--Nikodym cocycle is
$\widehat\Delta(\widehat u,\widehat v)=\Delta(\pi(\widehat u),\pi(\widehat v))$.
Pull back the partial coloring, graph, and matrix data. This includes
adjoining independent Bernoulli coordinates that the current coloring does
not use. Every bounded-cut ball calculation is then exactly a pullback.
Since $\pi$ preserves the base probability, the corresponding band
$L^1$ error norms agree. Uniqueness of the limits in
Lemma~\ref{lem:homogeneous-costs} therefore gives the same pullback identity
almost everywhere for each normalized limit field.

For a permutation $\sigma$ of the color names, extended by $\sigma(0)=0$,
the finite polynomial laws are relabeled by the same permutation. The
bounded-cut costs and their unique limits consequently satisfy
\[
 \bar m(\sigma\circ c)=\bar m(c),\qquad
 \bar d(\sigma\circ c)=\bar d(c),\qquad
 \bar d^{i}(\sigma\circ c)=\bar d^{\sigma^{-1}(i)}(c)
 \quad(1\leq i\leq r)
\]
almost everywhere. The first two identities express invariance of the
cost, while the last records how the virtual color options are renamed.

Define the nonnegative total cost of a partial coloring on the probability
base by
\begin{equation}
 \mathcal S(c)=\int_X\bar m(v)\,d\mu(v).
 \label{eq:coloring-total-cost}
\end{equation}
The preceding construction gives a uniform bound for this quantity.
The next lemma records the two approximation properties that will allow
us to use it: convergence after averaging over a long height strip, and
small integrated changes when two colorings agree on a large ball.

\begin{lemma}[Strip averages and locality]
\label{lem:strip-locality}
Under the assumptions of Lemma~\ref{lem:homogeneous-costs}, put
\[
 \Sigma_n=\{e^{-n}\leq W\leq e^n\},
 \qquad f_{\Sigma_n}=f^{e^{-n},e^n}.
\]
For every field $f$ in the list $m,d,d^1,\ldots,d^r$,
\begin{equation}
 \frac1{2n}\int_{\Sigma_n}
       |f_{\Sigma_n}-W\bar f|\,d\nu\longrightarrow0
 \quad(n\to\infty),
 \label{eq:coloring-strip}
\end{equation}
uniformly over the configurations and probability bases with the fixed
parameters.

If $c_1,c_2$ are two partial colorings on the same base with the same graph,
matrix, and cocycle data, and $E_R$ is the set where they agree on the
radius-$R$ graph ball, then
\begin{equation}
 \int_{E_R}|\bar f(c_1)-\bar f(c_2)|\,d\mu
      \leq\gamma(R),
 \qquad \gamma(R)\longrightarrow0,
 \label{eq:coloring-locality}
\end{equation}
where the same function $\gamma$ may be used for the finite list of fields
and is uniform over these configurations.

Both assertions hold after taking a product with an auxiliary probability
space and keeping the auxiliary parameter fixed along the equivalence
relation.  The cost fields may be chosen measurably on such product spaces.
\end{lemma}

\begin{proof}
In the strip integral there is the identity
\[
 |f_{\Sigma_n}-W\bar f|\,d\nu
    =|f_{\Sigma_n}/W-\bar f|\,d\mu\,ds.
\]
Divide the height interval $[-n,n]$ into unit intervals.  Translating each
to a fixed interval uses the homogeneity proved in
Lemma~\ref{lem:homogeneous-costs}.  For intervals a sufficiently large
fixed distance from both strip ends, the translated lower and upper
cutoffs are uniformly close to zero and infinity, respectively.  Uniform
joint convergence on the fixed band bounds their average errors by any
prescribed positive number.  The remaining bounded number of end
intervals has uniformly bounded error, and its contribution divided by
$2n$ tends to zero.  This proves Equation~\eqref{eq:coloring-strip}.

For Equation~\eqref{eq:coloring-locality}, first approximate both limiting
normalized fields on a fixed unit height band in $L^1$ by common, fixed
lower and upper cutoffs.  The errors are uniform by
Lemma~\ref{lem:homogeneous-costs}.  On this bounded cut,
Lemma~\ref{lem:bounded-costs} approximates each field uniformly by a
finite-radius calculation.  Agreement of the base colorings on a
sufficiently large ball makes these lifted local data identical.  Since
the normalized limit fields are independent of height, integration on
the unit band bounds the required base integral.  Choosing the cutoff
and then the radius for each prescribed accuracy gives
Equation~\eqref{eq:coloring-locality}.

Finally, all bounded-cut calculations are measurable, being limits of
rooted finite calculations.  They remain so with additional measurable
parameters.  An auxiliary parameter held fixed along the relation does
not alter the degree, edge ratios, or any spectral estimate.  The
preceding uniform estimates therefore hold on its product probability
space as well, and the $L^1$ limit construction gives jointly measurable
limit fields there.
\end{proof}

\subsection{The cost detects balls with excessive norm}

We next relate $\mathcal S$ to the norm conclusion in
Theorem~\ref{thm:bernoulli-coloring}.  The all-unassigned configuration has
zero cost.  Once a ball is fully assigned, a violation of the desired norm
bound forces a definite positive-root cost in that ball.  The invariant
lift allows us to sum this contribution without assuming that the
original probability measure is invariant.

\begin{lemma}[Initial cost and bad balls]
\label{lem:bad-balls}
Assume the hypotheses and parameter choices of
Theorem~\ref{thm:bernoulli-coloring}.  The all-unassigned partial coloring
has $\mathcal S=0$.  Fix $m\geq1$ and $\alpha>0$.  For any measurable
partial coloring $c$, let $\operatorname{Bad}_m(c)$ be the set of vertices
$v$ for which every vertex of $B_G(v,m)$ is assigned and
\[
 \|T_{B_G(v,m),c}\|>2\beta t+\alpha.
\]
There is a finite constant $C_4$, depending only on $m,D_0,K,\alpha$, such
that
\begin{equation}
 \mu(\operatorname{Bad}_m(c))\leq C_4\mathcal S(c).
 \label{eq:coloring-bad-balls}
\end{equation}
The same assertions hold on probability extensions with the pulled-back
graph, matrix, and cocycle data.
\end{lemma}

\begin{proof}
For the all-unassigned coloring, Lemma~\ref{lem:finite-cost-endpoints}
gives zero cost on every finite principal restriction. Finite restriction
limits and Lemma~\ref{lem:homogeneous-costs} therefore give $\mathcal S=0$.

To prove Equation~\eqref{eq:coloring-bad-balls}, partition the bad centers
into a bounded finite number of measurable classes so that two distinct
centers in one class have graph distance greater than $2m+1$.  Here is a
measurable construction.  A bounded-degree Borel graph has a countable
measurable partition into independent sets: assign to each vertex a
finite prefix of an injective standard-space code long enough to
separate it from all of its neighbors, and include the prefix length in
the assigned code.  Equal assigned codes cannot occur at adjacent
vertices.  Visit these independent sets in sequence, and give each
vertex the first color not already used by one of its previously colored
neighbors, from a palette of one more than the degree bound.  This is a
measurable proper coloring.  Apply this construction to the graph whose
edges join vertices at distance at most $2m+1$, and intersect its color
classes with $\operatorname{Bad}_m(c)$.  The number of resulting classes
is bounded in terms of $D_0,m$.

Fix one such class $C$ of bad centers.  In the invariant lift retain the
strip $\Sigma_n$, and select the lifted centers above $C$ whose height is
at distance at least $m\log K$ from both ends of $[-n,n]$.  Every
radius-$m$ ball around a selected center lies entirely in the strip.
Moreover these balls are disjoint and nonadjacent along each lifted
orbit.  Indeed Equation~\eqref{eq:coloring-height-change} and the cocycle
law give exactly one lifted site above each base site in a fixed lifted
orbit.  Distinct centers there therefore project to distinct centers of
$C$, whose distance is greater than $2m+1$.

Delete every strip site outside these balls.  By
Lemma~\ref{lem:integral-change}, deletion cannot increase the total
integrated cost, so
\begin{equation}
 \int_{\Sigma_n}m_{\Sigma_n}\,d\nu
       \geq\int m_{\mathrm{retained}}\,d\nu.
 \label{eq:coloring-ball-deletion}
\end{equation}
Each retained graph component is one of the selected finite balls.
Invariant mass transport sends the rooted costs in each such ball to
its center, and hence the right side of
Equation~\eqref{eq:coloring-ball-deletion} equals the integral over the
selected lifted centers of $s(P_{\mathrm{ball}})$.

For a bad ball, the second part of Lemma~\ref{lem:finite-cost-endpoints}
gives $s(P_{\mathrm{ball}})\geq(\alpha/2)\min_{\mathrm{ball}}W$.
The edge-ratio bound then gives
\begin{equation}
 s(P_{\mathrm{ball}})
    \geq\frac\alpha2 K^{-m}W_{\mathrm{center}}.
 \label{eq:coloring-ball-root-cost}
\end{equation}

Divide Equation~\eqref{eq:coloring-ball-deletion} by $2n$ and let
$n\to\infty$.  Equation~\eqref{eq:coloring-strip} makes the left side
tend to $\mathcal S(c)$.  On the right,
$W\,d\nu=d\mu\,ds$, and the allowed height interval has length
$2n-2m\log K$ for all sufficiently large $n$.  Thus
Equation~\eqref{eq:coloring-ball-root-cost} gives
\[
 \frac\alpha2 K^{-m}\mu(C)\leq\mathcal S(c).
\]
Summing over the bounded number of classes proves
Equation~\eqref{eq:coloring-bad-balls}.  Every part of the argument uses
only the stated probability-space, cocycle, degree, and matrix bounds,
so it applies unchanged to the indicated extensions.
\end{proof}

\subsection{Assigning a small fraction of the remaining vertices}

We now use Equation~\eqref{eq:coloring-limit-assignment} to choose colors.
Changing one site has a controlled cost, but simultaneous changes must be
compared in their intermediate environments.  Independent uniform times
provide an order for those comparisons.  We first perform the comparison
on a bounded height strip, where the costs have deterministic definitions
by finite-neighborhood limits.  Only then do we pass to the homogeneous
costs.  This order of operations will avoid evaluating a version of an
$L^1$ limit at a random value of its parameter.

Fix a current partial coloring $c$ determined by coordinates already used,
and let $E_0=\{c=0\}$.  At every site in $E_0$, choose a color $i(v)$
minimizing $\bar d^{i}(v)$.  Break ties by a fresh uniform random
permutation of the colors, independently at each site, taking the first
minimizer in that permutation.  The costs of $c$ are first defined on the
space carrying the old coordinates and then pulled back to the space
carrying the extra tie coordinates.  The bounded-cut definitions and
their limits commute with this pullback.  In particular, the choices
$i(v)$ use only old and tie coordinates.

Now add fresh independent uniform times $U_v\in(0,1)$ at the sites,
independent of both the old and the tie coordinates.  For $0\leq\theta<1$,
define $c^\theta$ by assigning $i(v)$ at each site $v\in E_0$ with
$U_v<\theta$ and leaving all other sites as in $c$.  All these definitions
are measurable on the relative Bernoulli extension.  Null exceptions can
be filled by measurable default values before adding the next independent
coordinates; their saturations are negligible.  We use this convention
throughout the finite sequence of steps below.

\begin{lemma}[One assignment step]
\label{lem:hazard-step}
Under the hypotheses of Theorem~\ref{thm:bernoulli-coloring}, let $c$ be
any partial coloring determined before the fresh times, and choose
$i(v)$ and $c^\theta$ as just described.  There is a nonnegative function
$g$, depending only on the fixed graph and polynomial parameters, such
that $g(h)\to0$ as $h\downarrow0$ and, for every $0<h<1/2$,
\begin{equation}
 \mathcal S(c^h)
      \leq(1+24h)\mathcal S(c)+h g(h).
 \label{eq:coloring-step-bound}
\end{equation}
The same function $g$ works for all such steps and current configurations.
\end{lemma}

\begin{proof}
Apply Lemma~\ref{lem:integral-change} to the change from $c$ to $c^h$ on
the invariant lift restricted to $\Sigma_n$.  Order the changed sites by
their times $U_v$.  These times are almost surely distinct along every
countable orbit, after discarding a saturated null set.  The order need
not be a well-order, as allowed in that lemma.  At a comparison site $v$
the pre-change configuration is the restriction of $c^{U_v}$, with the
number $U_v$ held fixed as a threshold along that orbit.

We claim that the resulting integral of the single-site increments is
\begin{equation}
 \begin{split}
 &\int_{\Sigma_n}
       \bigl(m_{\Sigma_n}(c^h)-m_{\Sigma_n}(c)\bigr)\,d\nu\\
 &\quad=\int_0^h\frac{d\theta}{1-\theta}
       \int_{\Sigma_n}1_{\{v\in E_0,\ U_v\geq\theta\}}
       \left[
        d_{\Sigma_n}^{i(v)}(c^\theta)
                         -d_{\Sigma_n}(c^\theta)
       \right]d\nu.
 \end{split}
 \label{eq:coloring-finite-strip-step}
\end{equation}
Here $\nu$ includes all label laws as well as the height measure, and $v$
denotes the current base site of a lifted root.

To justify the claim, condition at a root in $E_0$ on the old coordinates, all tie
coordinates, and every fresh time except the time at that root.  Hold
the height fixed as well.  For a fixed threshold $\theta$, leave the root
site unassigned, and denote the bracketed difference in
Equation~\eqref{eq:coloring-finite-strip-step} by $D(\theta)$ for these
fixed other data.  Whenever the root time $u_0$ is at least $\theta$,
this difference is independent of its actual value.  The chosen colors
were determined before the fresh times, other sites use only their own
fresh times in the threshold test, and neither the graph nor the
restricted weight data use this new root coordinate.  This independence
holds for the deterministic rooted finite-subset calculations and hence
for the bounded-strip costs obtained from them.  The same description
applies at $u_0=\theta$ because the update condition is strict.

When $u_0<h$, the ordered comparison increment at the root is therefore
$D(u_0)$, outside the null event of a tie with another site.  Integrating
over this last uniform time uses precisely the identity
\[
 \int_0^h D(u_0)\,du_0
   =\int_0^h\frac{d\theta}{1-\theta}
                     \int_\theta^1 D(\theta)\,du_0.
\]
The bounded-strip computations are measurable jointly in the threshold
and in the remaining data, again by their finite rooted definitions.
Fubini's Theorem, first for almost every root and height and then over
the remaining coordinates, proves
Equation~\eqref{eq:coloring-finite-strip-step}.  This argument takes place
entirely before choosing any homogeneous limit version depending on
$\theta$.

We can now pass to that limit in an integrated expression.  Divide
Equation~\eqref{eq:coloring-finite-strip-step} by $2n$ and use
Lemma~\ref{lem:strip-locality}.  For the right side, adjoin $\theta$ with
normalized Lebesgue probability on $[0,h]$, holding it fixed along the
equivalence relation, and regard $c^\theta$ as one configuration on this
augmented base.  The lemma applies uniformly there.  Its absolute
integrated errors also control the errors after multiplication by the
displayed mask, after selection among the finite list of options
$i(v)$, and after multiplication by $(1-\theta)^{-1}\leq2$.
For a height-independent limit field, cancellation of $W$ against the
height density leaves a height interval of length exactly $2n$.
Consequently
\begin{equation}
 \begin{split}
 \mathcal S(c^h)-\mathcal S(c)
    =\int_0^h\frac{d\theta}{1-\theta}
      \int_X1_{\{v\in E_0,\ U_v\geq\theta\}}
       \left[
        \bar d^{i(v)}(c^\theta)-\bar d(c^\theta)
       \right]d\mu.
 \end{split}
 \label{eq:coloring-limit-step}
\end{equation}
All fields in this formula have jointly measurable versions on the
augmented space.  No evaluation of an unspecified limit at the random
diagonal $\theta=U_v$ has been used.

It remains to compare the intermediate costs with the old costs, for
which $i(v)$ was chosen.  Let $D_R$ be a finite bound, depending only on
$D_0,R$, for the cardinality of a radius-$R$ graph ball.  On the event
that this ball has no update from $c$ to $c^\theta$, the two partial
colorings agree there.  Equation~\eqref{eq:coloring-locality} controls the
integrated difference of their cost fields on that event; summing over
the finite option list also controls the selected field.  Conditional
on the old and tie coordinates, the probability of an update anywhere
in the ball is at most $D_R h$, by the union bound for the fresh uniform
times.  On the event that an update does occur, the cost differences are uniformly bounded
by Lemma~\ref{lem:homogeneous-costs}.

These observations apply either for each fixed $\theta$, or jointly
with normalized measure on $[0,h]$ as above.  They show that replacing
the two costs in the brackets of Equation~\eqref{eq:coloring-limit-step}
by their values at $c$ creates an error in absolute value at most
\begin{equation}
 C h\bigl(\gamma(R)+D_Rh\bigr),
 \label{eq:coloring-step-error}
\end{equation}
where $C$ depends only on the fixed parameters.  The factor
$(1-\theta)^{-1}$ is bounded by $2$.

The old costs and the minimizing choice $i(v)$ are independent of the
fresh root time.  The survival probability $1-\theta$ therefore cancels
the denominator in Equation~\eqref{eq:coloring-limit-step}.  By
Equation~\eqref{eq:coloring-limit-assignment}, the resulting old-cost
contribution is at most
\[
 h\int_{E_0}\bigl(\bar d^{i(v)}(c)-\bar d(c)\bigr)\,d\mu(v)
       \leq24h\int_{E_0}\bar m(c)\,d\mu
       \leq24h\mathcal S(c).
\]
For each prescribed accuracy in Equation~\eqref{eq:coloring-step-error},
first take $R$ large enough to control $\gamma(R)$ and then take $h$
small enough to control $D_Rh$.  Equivalently, taking the infimum of
these uniform error bounds over integer $R$ gives a nonnegative function
$g(h)\to0$.  This proves Equation~\eqref{eq:coloring-step-bound} uniformly
over the current configurations.
\end{proof}

\subsection{Completion of the coloring}

The estimates now have the two features needed for a finite construction.
Repeated small assignment steps can cover any prescribed ball with high
probability, and, for a fixed total assignment time, their accumulated
cost tends to zero with the step size.  The bad-ball estimate then supplies
the norm conclusion.

\begin{proof}[Proof of Theorem~\ref{thm:bernoulli-coloring}]
Work on the relative Bernoulli extension of the original base, reserving
independent coordinates at every site for each step.  Start with the
all-unassigned coloring, whose cost is zero by Lemma~\ref{lem:bad-balls}.
Choose $T_0>0$ and $h\in(0,1/2)$, and apply the construction of
Lemma~\ref{lem:hazard-step} with independent fresh coordinates for
\[
 N=\lceil T_0/h\rceil
\]
steps.  Write $S_j$ for the cost after $j$ steps.  Iterating
Equation~\eqref{eq:coloring-step-bound} gives
\begin{equation}
 S_N\leq\frac{g(h)}{24}\bigl((1+24h)^N-1\bigr).
 \label{eq:coloring-iteration}
\end{equation}
For fixed $T_0$ this tends to zero as $h\downarrow0$, since
$(1+24h)^N\leq\exp(24(T_0+h))$.

An unassigned site receives a color at its first successful hazard.
Conditional on the original base point, its probability of remaining
unassigned after all $N$ steps is $(1-h)^N\leq e^{-T_0}$.  This conclusion
does not depend on the color selected at a successful step.  If $D_m$
bounds the size of a radius-$m$ ball, the conditional probability that
any site in that ball remains unassigned is at most $D_m e^{-T_0}$.
The unconditional probability that the ball is fully assigned but
violates the required norm bound is at most $C_4 S_N$, by
Lemma~\ref{lem:bad-balls}.  First choose $T_0$ large, and then $h$ small,
so that
\[
 D_m e^{-T_0}+C_4S_N<\eta.
\]
Complete the partial coloring by fresh independent uniform colors at the
remaining sites.  Every previously complete ball is unchanged.  The
resulting coloring therefore satisfies
Equation~\eqref{eq:coloring-ball-conclusion}.

It remains to verify the conditional marginals. Fix a permutation
$\sigma$ of $\{1,\ldots,r\}$, and let $\Theta_\sigma$ apply it to the color
names in every tie-breaking permutation and every final uniform color input,
leaving all times unchanged. This transformation fixes the original base
point, commutes with rerooting, and preserves each conditional product measure $\kappa_v$. On the
height lift it also fixes the height, so it acts isometrically on every
band $L^1$ space.

Suppose inductively that the current coloring satisfies
$c\circ\Theta_\sigma=\sigma\circ c$. Exact pullback of the finite-cut
calculations and their color covariance give
\[
 \bigl((d^{\sigma(i)})^{l,H}(c)\bigr)\circ\Theta_\sigma
       =(d^i)^{l,H}(c),\qquad 1\leq i\leq r.
\]
The band $L^1$ isometry and uniqueness of the homogeneous limits give
$(\bar d^{\sigma(i)}(c))\circ\Theta_\sigma=\bar d^i(c)$ almost everywhere.
Thus the set of minimizing colors at the transformed input is the
$\sigma$-image of the original minimizing set. The transformed
permutation tie-breaker chooses the corresponding color, and the update
test is unchanged because the times are unchanged.

There are only finitely many color permutations and, with $T_0,h$ fixed,
finitely many steps. Choose a common conull set for the required identities,
intersect its images under the finite color action, and discard the saturation
of its complement under the relation. The induction just described then
holds throughout the construction on that set, starting from the
all-unassigned coloring. The final uniform color inputs transform in the
same way, so the completed coloring also obeys
$c\circ\Theta_\sigma=\sigma\circ c$.

Since $\Theta_\sigma$ preserves the conditional product measure over each
original base point, disintegration shows that all output colors at the
root have equal conditional probabilities for almost every base point.
They sum to one, which proves Equation~\eqref{eq:coloring-marginal} and
completes the proof.
\end{proof}
